\section{Standard Model Descent, Generations, and Fermion Masses}
\label{sec:sm:main}

This section develops the Standard Model descent within UCFT. We start
from the vacuum-selection theorem (\autoref{sec:potential:main}), which fixes
the residual gauge group $H = \mathrm{Spin}(10)\times U(1)$ and the matter
kinematics, and we spell out the symmetry-breaking cascade down to
$SU(3)_c\times U(1)_{\mathrm{em}}$. The descent chain, the branchings, and the
anomaly arithmetic are \emph{theorems} given the residual group and the assigned
charges; chirality (P4), the three-generation horizontal symmetry (P5), and the
GUT--Higgs/Yukawa sector (P6) are \emph{added postulates}; the heterotic
generation match and the listed phenomenology are \emph{conjectural
correspondences}. \autoref{sec:sm:summary} summarizes the epistemic status.
Conventions are fixed globally: metric signature $(-,+,+,+)$, $[f]=$ mass, the
$32$ coset Goldstones are the \emph{clock fields}, and $M_{\mathrm{GUT}}$ denotes
the breaking-VEV scale tied to the decay constant $f$.

\subsection{The explicit symmetry-breaking chain}
\label{sec:sm:descent}

\begin{postulate}[GUT--Higgs sector, P6 part I]\label{post:sm:P6higgs}
The descent below $H=\mathrm{Spin}(10)\times U(1)$ proceeds through the real
$SO(10)$ Higgs multiplets $45_H$ \emph{or} $54_H$, $126_H$, and $10_H$, whose VEVs
realize the cascade of \autoref{thm:sm:descent}. These multiplets are added
structure, not consequences of the two axioms.
\end{postulate}

\begin{theorem}[Descent chain]\label{thm:sm:descent}
With the residual group $H=\mathrm{Spin}(10)\times U(1)$ fixed by the proved
rank-$1$ idempotent vacuum $M=E_6/(\mathrm{Spin}(10)\times U(1))=\mathrm{EIII}$
(\autoref{thm:potential:vacuum}), and granting the Higgs content of
\autoref{post:sm:P6higgs}, the gauge symmetry descends as
\begin{equation}\label{eq:sm:chain}
\begin{aligned}
E_6
 &\;\xrightarrow{\;\text{P1: rank-}1\text{ idempotent vacuum }M=\mathrm{EIII}\;}\;
 \mathrm{Spin}(10)\times U(1) \\
 &\;\xrightarrow{\;45_H\ \text{or}\ 54_H\;}\;
 SU(5)\times U(1)_X \\
 &\;\xrightarrow{\;126_H\ (\text{breaks }B\!-\!L\text{ and rank})\;}\;
 SU(3)_c\times SU(2)_L\times U(1)_Y \\
 &\;\xrightarrow{\;10_H\ (\text{electroweak},\ v=246\ \mathrm{GeV})\;}\;
 SU(3)_c\times U(1)_{\mathrm{em}}\,.
\end{aligned}
\end{equation}
\end{theorem}

\begin{proof}[Proof sketch]
Step $0$ ($E_6\to \mathrm{Spin}(10)\times U(1)$) is not a Higgs VEV in the GUT sense
but the vacuum-selection step: the $E_6$-invariant potential $V$ of P1 has its global
minimum on the rank-$1$ primitive-idempotent orbit, whose stabilizer is exactly
$H=\mathrm{Spin}(10)\times U(1)$ (\autoref{thm:potential:vacuum}). Steps $1$--$3$
are standard $SO(10)$ symmetry breaking by real Higgs reps: $\langle 45_H\rangle$ (or
$\langle 54_H\rangle$) reduces $SO(10)$ to its $SU(5)$ subgroup; the SM-singlet,
$B\!-\!L=2$ component of $126_H$ acquires a VEV $v_R\sim M_{\mathrm{GUT}}$, breaking
$B\!-\!L$ and lowering the rank by one (removing the extra Abelian factor and supplying
$M_R$, \autoref{thm:sm:seesaw}); $\langle 10_H\rangle=v=246$~GeV is ordinary
electroweak breaking. Each embedding is a standard maximal/regular subgroup chain;
no step uses a forbidden $\{16,16bar\}$ Higgs.
\end{proof}

\begin{theorem}[Hypercharge embedding and GUT normalization]\label{thm:sm:hypercharge}
Hypercharge is the standard $SU(5)$-embedded Cartan combination. Acting on the
$SU(5)$ fundamental $5$,
\begin{equation}\label{eq:sm:TY}
T_Y \;=\; \tfrac{1}{\sqrt{60}}\,\mathrm{diag}(-2,-2,-2,\,3,\,3)\,,
\qquad
Y \;=\; \mathrm{diag}\!\big(-\tfrac13,-\tfrac13,-\tfrac13,\,\tfrac12,\tfrac12\big)
\end{equation}
on $5bar=(d^c,l)$. In the canonical group-theoretic (GUT) normalization the unified
$U(1)$ coupling obeys
\begin{equation}\label{eq:sm:alpha1}
\alpha_1 \;=\; \tfrac{5}{3}\,\alpha_Y\,,
\qquad\Longleftrightarrow\qquad
g_1^2 \;=\; \tfrac{5}{3}\,g_Y^2 .
\end{equation}
\end{theorem}

The orthogonal $U(1)_X$ (the $SO(10)/SU(5)$ Abelian factor, orthogonal to $Y$) is
broken at the $126_H$ scale as a heavy $Z'$, or kept light only if separately
anomaly-checked. A scale remark that we use repeatedly downstream: $M_{\mathrm{GUT}}$
is the \emph{breaking-VEV scale tied to $f$}, \emph{not} the imported MSSM
$2\times10^{16}$~GeV crossing. Below $M_{\mathrm{GUT}}$ one runs the \emph{three}
Standard-Model $\beta$-functions, not a single $b_0$.

\subsection{One family from the \texorpdfstring{$16$}{16}}
\label{sec:sm:family}

\begin{theorem}[$16$ = one SM family $+$ $\nu^c$]\label{thm:sm:family}
Under $SO(10)\to SU(5)\times U(1)_X$ the chiral spinor decomposes as
\begin{equation}\label{eq:sm:16branch}
16 \;\longrightarrow\; 10_{+1}\;+\;\overline{5}_{-3}\;+\;1_{+5}\,,
\end{equation}
and further under $SU(5)\to SU(3)_c\times SU(2)_L\times U(1)_Y$ (with $Q=T_3+Y$) the
content is exactly one Standard-Model family plus a singlet:
\begin{equation}\label{eq:sm:family}
\begin{aligned}
10 &\;\supset\; q=(3,2)_{+1/6},\quad u^c=(\bar3,1)_{-2/3},\quad e^c=(1,1)_{+1},\\
\overline{5} &\;\supset\; d^c=(\bar3,1)_{+1/3},\quad l=(1,2)_{-1/2},\\
1 &\;=\; \nu^c=(1,1)_{0}\ \ (\text{right-handed neutrino, $SU(5)$ singlet}).
\end{aligned}
\end{equation}
Thus $10+\overline{5}+1=16$ Weyl fermions $=(q,u^c,d^c,l,e^c)+\nu^c$: one SM family
plus a right-handed neutrino, with no mirror partners.
\end{theorem}

The chiral family sits inside $16_{+1}$ of the matter $27$ of $E_6$, whose canonical
UCFT branching (\autoref{sec:potential:branchings}) is
\begin{equation}\label{eq:sm:27branch}
27 \;=\; 1_{+4}\;+\;10_{-2}\;+\;16_{+1}\,,
\end{equation}
where the $1_{+4}$ is an extra SM-singlet modulus available as a right-handed-neutrino
modulus. That a single $16$ rather than a vector-like $16\oplus\overline{16}$ appears
per family is the content of the chirality postulate:

\begin{postulate}[Matter content and chirality, P4]\label{post:sm:P4}
Each family is one chiral $16$ of $\mathrm{Spin}(10)$ with the $U(1)$ charge of
\eqref{eq:sm:27branch}; there are no mirror partners.
\end{postulate}

\subsection{Three generations: \texorpdfstring{$SU(3)_F$}{SU(3)F} on
\texorpdfstring{$J_3(\mathbb{O})\otimes\mathbb{C}^3$}{J3(O)xC3}}
\label{sec:sm:generations}

\begin{postulate}[Family symmetry, P5]\label{post:sm:P5}
The matter carries an $SU(3)_F$ horizontal (family) symmetry acting on
$J_3(\mathbb{O})\otimes\mathbb{C}^3$; the three copies in the $\mathbb{C}^3$ factor are
three chiral $16$s of $\mathrm{Spin}(10)$, with no mirror partners. Three generations
are \emph{not} a consequence of the two axioms; they are added structure.
\end{postulate}

\begin{theorem}[Recomputed $n_F$ and three-generation $b_0$]\label{thm:sm:b0}
With one chiral $\mathbf{16}$ per family (P4, P5), the fermion count $n_F$ in the
one-loop $\mathrm{Spin}(10)$ $\beta$-function is the number of chiral
$\mathbf{16}$ \emph{representations} (one per generation), not the Weyl-fermion
dimension:
\begin{equation}\label{eq:sm:nF}
n_F \;=\; \#\text{generations} \;=\; 3
\qquad\text{(three chiral $\mathbf{16}$s, each weighted at $T(\mathbf{16})=2$).}
\end{equation}
The gauge $\beta$-function (canonical $+2$ term restored; asymptotically free, no
interacting fixed point) is
\begin{equation}\label{eq:sm:beta}
\beta_{\hat g^{2}} \;=\; 2\,\hat g^{2} \;+\; \frac{b_0}{48\pi^2}\,\hat g^{4},
\qquad
b_0 \;=\; \tfrac{11}{3}C_A \;-\; \tfrac{4}{3}T\,n_F \;-\; \tfrac{1}{6}T\,n_S\,,
\end{equation}
with the locked Casimirs $C_A(\mathrm{Spin}(10))=8$, $T(\mathbf{16})=2$,
$C_2(\mathbf{16})=45/8$. At the GUT matching scale with the full UCFT scalar
content, $n_S=28$ massive coset scalars enter the loop. Direct evaluation gives
\begin{equation}\label{eq:sm:b0three}
\boxed{\,
b_0(3\,\mathrm{gen})
=\tfrac{11}{3}(8)-\tfrac{4}{3}(2)(3)-\tfrac{1}{6}(2)(28)
=\tfrac{88-24-28}{3}=12\,.}
\end{equation}
The canonical one-generation anchor $b_0(1\,\mathrm{gen})=26$ uses $n_F=1$ and
$n_S=2$ (minimal GUT Higgs/Yukawa scalars at that scale):
\begin{equation}\label{eq:sm:b0one}
b_0(1\,\mathrm{gen})
=\tfrac{11}{3}(8)-\tfrac{4}{3}(2)(1)-\tfrac{1}{6}(2)(2)=\tfrac{78}{3}=26.
\end{equation}
We use $b_0=12$ for three-generation running below $M_{\mathrm{GUT}}$ and in the
AS stability check of \autoref{sec:gravity:stability}.
\end{theorem}

\begin{remark}
The gauge sector is asymptotically \emph{free}: the canonical $+2\,\hat g^{2}$
term in \eqref{eq:sm:beta} drives a Gaussian UV point $\hat g^{2}_\star=0$; there
is no interacting gauge fixed point. More fermions \emph{decrease} $b_0$ in
\eqref{eq:sm:beta}, so $b_0(3)<b_0(1)$. UV completeness is asymptotic freedom
plus the (indicative) gravitational fixed point $\kappa_\star=96\pi^2/5$
(\autoref{sec:gravity:AS}) plus irrelevant higher operators.
\end{remark}

\begin{postulate}[Hierarchy and mixing from flavon VEVs, P5 part II]\label{post:sm:flavons}
$SU(3)_F$ is broken sequentially ($SU(3)_F\to SU(2)_F\to \mathbf{1}$) by flavon VEVs,
generating the inter-generation Yukawa hierarchy and the CKM (quark) and PMNS (lepton)
mixing matrices as ratios of flavon VEVs to the family scale.
\end{postulate}

\begin{indicative}[Heterotic corroboration]\label{ind:sm:heterotic}
The heterotic $\mathbb{Z}_5$ quintic $\hat Q=Q/\mathbb{Z}_5$ compactification with the
tangent-bundle $SU(3)$ embedding gives the same generation count,
\begin{equation}\label{eq:sm:het}
h^1(\hat Q,V_{16})=3,\qquad h^1(\hat Q,V_{10})=14,\qquad h^1(\hat Q,\mathbf{1})=1,
\end{equation}
with the single bundle singlet an overall right-handed-neutrino modulus and the
$4+28=32$ heterotic moduli matching the $4$ clock-field Goldstone $+$ $28$ massive
coset scalars of UCFT. The low-energy central charges $(c_L,c_R)=(16,10)$ enter here.
This is a correspondence, not a duality, and not a derivation of the generation number.
\end{indicative}

\subsection{Yukawa and seesaw sector (P6)}
\label{sec:sm:yukawa}

\begin{theorem}[No $\{16,16bar\}$ Yukawa]\label{thm:sm:noyukawa}
The naive Yukawa $-y\,\bar\psi\,\phi^a T_a\,\psi$ with $\phi\in\{16,16bar\}$ is
forbidden: it is non-$H$-invariant and gives zero mass, because
\begin{equation}\label{eq:sm:tensor}
16\times 16 \;=\; 10+120+126\,,
\qquad
16\times \overline{16} \;=\; 1+45+210\,,
\end{equation}
so a $\{16,16bar\}$ scalar furnishes no $H$-singlet, and $H=\mathrm{Spin}(10)\times
U(1)$ is unbroken at the vacuum. The viable couplings use the real GUT-Higgs reps.
\end{theorem}

\begin{postulate}[Yukawa couplings, P6 part II]\label{post:sm:P6yukawa}
The family-indexed Yukawa Lagrangian is
\begin{align}
\mathcal{L}^{\mathrm{Dirac}}_{\mathrm{Yuk}}
 &\;=\; y^{(10)}_{ij}\,16_i\,16_j\,10_H \;+\; \mathrm{h.c.}\,,
 \label{eq:sm:Ldirac}\\
\mathcal{L}^{\mathrm{Maj}}_{\mathrm{Yuk}}
 &\;=\; y^{(126)}_{ij}\,16_i\,16_j\,\overline{126}_H \;+\; \mathrm{h.c.}
 \;\;\;\Big(\text{or}\;\; \tfrac{1}{M_*}\,16_i\,16_j\,\overline{16}_H\,\overline{16}_H\Big)\,,
 \label{eq:sm:Lmaj}
\end{align}
both $SO(10)$-invariant since $16\times 16\supset 10$ and $16\times 16\supset 126$,
with $i,j$ the $SU(3)_F$ family indices.
\end{postulate}

\begin{theorem}[Type-I seesaw]\label{thm:sm:seesaw}
After electroweak breaking $\langle 10_H\rangle=v=246$~GeV the Dirac masses are
$m_D\sim y^{(10)}v$. The SM-singlet, $B\!-\!L=2$ component of $126_H$ takes
$\langle\,\cdot\,\rangle=v_R\sim M_{\mathrm{GUT}}$ (step~2 of \eqref{eq:sm:chain}),
giving the right-handed Majorana mass $M_R\sim y^{(126)}v_R$ (equivalently
$M_R\sim v_R^2/M_*$ from the dim-$5$ operator). The $(\nu,\nu^c)$ mass matrix
\begin{equation}\label{eq:sm:matrix}
\mathcal{M}_\nu \;=\;
\begin{pmatrix} 0 & m_D \\ m_D^{\mathsf T} & M_R \end{pmatrix}
\end{equation}
yields, to leading order in $m_D/M_R$, the light-neutrino mass
\begin{equation}\label{eq:sm:mnu}
m_\nu \;\simeq\; \frac{m_D^2}{M_R} \;\sim\; 0.05\ \mathrm{eV}\,,
\end{equation}
the observed atmospheric scale for $m_D\sim v$ and $M_R\sim 10^{14}$--$10^{15}$~GeV.
PMNS mixing follows from the flavon structure (\autoref{post:sm:flavons}).
\end{theorem}

These multiplets are added physics; stating them is what makes the construction an
embedding rather than a derivation.

\subsection{Anomaly arithmetic}
\label{sec:sm:anomaly}

\begin{theorem}[Anomaly cancellation]\label{thm:sm:anomaly}
The composite $SO(10)\times U(1)$ gauge sector is anomaly-free. With Dynkin indices
$A_{16}=A_{10}=1$, $A_1=0$ and the $27$-charges of \eqref{eq:sm:27branch}, the mixed
$U(1)$--$SO(10)^2$ coefficient is
\begin{equation}\label{eq:sm:mixed}
(+1)A_{16}+(-2)A_{10}+(+4)A_1
 \;=\;(+1)(1)+(-2)(1)+(+4)(0)\;=\;-1\,,
\end{equation}
cancelled by the Green--Schwarz mechanism: the $E_6$-invariant $3$-form $\Omega_3$
pulls back to $dB=(1/2\pi)F_1$, and
$S_{\mathrm{GS}}=\int B\wedge \mathrm{Tr}(F\wedge F)/(8\pi^2)$ has a gauge variation
exactly cancelling the $-1$. The cubic $U(1)^3$ coefficient vanishes identically,
\begin{equation}\label{eq:sm:cubic}
16(+1)^3+10(-2)^3+1(+4)^3 \;=\; 16-80+64 \;=\; 0\,.
\end{equation}
Below $M_{\mathrm{GUT}}$ each family is the anomaly-free $16$ (the $\nu^c$ is exactly
what makes the SM hypercharge anomalies cancel family-by-family).
\end{theorem}

\begin{remark}[Strings and topology, correspondence]\label{rem:sm:strings}
The composite $SO(10)\times U(1)$ connection (Maurer--Cartan), the Green--Schwarz
cancellation, the vortex strings, and the heterotic worldsheet CFT with
$(c_L,c_R)=(16,10)$ are retained as a low-energy correspondence. Topology caveat:
$\mathrm{EIII}$ is simply connected, so $\pi_1(M)=\mathbb{Z}$ is \emph{false} of
$\mathrm{EIII}$ itself; the vortex strings arise from the gauged $U(1)$-quotient
configuration space (or are softened to a correspondence).
\end{remark}

\subsection{Phenomenology (indicative)}
\label{sec:sm:pheno}

\begin{indicative}[Proton decay, leptogenesis, monopole dilution]\label{ind:sm:pheno}
Once descent and seesaw exist, the leading observables are, all indicative:
\begin{equation}\label{eq:sm:proton}
\tau(p\to e^+\pi^0)\;\sim\;\frac{M_{\mathrm{GUT}}^4}{\alpha_{\mathrm{GUT}}^2\,m_p^5}\,,
\end{equation}
to be confronted with Super-/Hyper-Kamiokande, with the scalar-leptoquark channel
forbidden by matter parity / $B\!-\!L$; baryogenesis via leptogenesis from the heavy
$\nu^c$ decays (\autoref{thm:sm:seesaw}); and monopole dilution by inflation, with the
inflaton identified as one of the $28$ massive coset scalars
(\autoref{sec:gravity:spectrum}) and inflation occurring \emph{after} the GUT
transition. Here $M_{\mathrm{GUT}}$ is the breaking-VEV scale tied to $f$.
\end{indicative}

\subsection{Summary of results}
\label{sec:sm:summary}

\begin{itemize}
\item \textbf{Theorems (proved within UCFT):} the $E_6\to\mathrm{Spin}(10)\times U(1)$
vacuum selection (\autoref{thm:potential:vacuum}); the descent
(\autoref{thm:sm:descent}); the branchings of $27$ and $16$
(\autoref{thm:sm:family}); the anomaly arithmetic
(\autoref{thm:sm:anomaly}, mixed $=-1$ cancelled by GS, cubic $=0$); and the
three-generation $b_0=12$ (\autoref{thm:sm:b0}).
\item \textbf{Postulates (added structure):} P4 (one chiral $16$/family, no mirrors,
\autoref{post:sm:P4}); P5 ($SU(3)_F\to$ three generations, $n_F=3$,
\autoref{post:sm:P5}); P6 ($10_H$ Dirac, $126_H$/dim-$5$ Majorana, type-I seesaw,
\autoref{post:sm:P6higgs}, \autoref{post:sm:P6yukawa}).
\item \textbf{Correspondence (indicative):} heterotic $h^1(V_{16})=3$ and the
$4+28=32$ moduli match (\autoref{ind:sm:heterotic}); the phenomenology of
\autoref{ind:sm:pheno}.
\end{itemize}

\noindent\textbf{Conclusion.} Under postulates P4--P6, UCFT embeds the Standard
Model: three anomaly-free chiral families, standard hypercharge with GUT
normalization $\alpha_1=\tfrac53\alpha_Y$, electroweak breaking at $v=246$~GeV, and a
type-I seesaw $m_\nu\sim0.05$~eV.
